堆排序（Heapsort）把数组想象成一棵二叉树——
这就是「堆」这个名字的来源。它不需要额外的数据结构，
只在原数组上通过下标关系操作。

\subsection{核心思路：把数组看成一棵二叉树}

堆排序用一种巧妙的下标关系：

\begin{itemize}
    \item 节点 \texttt{i} 的\textbf{父节点}是 \texttt{i/2}
    \item 节点 \texttt{i} 的\textbf{左孩子}是 \texttt{2*i}
    \item 节点 \texttt{i} 的\textbf{右孩子}是 \texttt{2*i+1}
\end{itemize}

这个映射不需要额外内存，只靠四则运算。
注意它沿用了《算法导论》的\textbf{从 1 开始}的数组习惯——
\texttt{A[0]} 通常放一个哨兵或被忽略。

一个\textbf{最大堆}满足：每个父节点 \verb|>=| 它的两个孩子，
即 \texttt{A[Parent(i)] >= A[i]} 对所有 \texttt{i} 成立。
这保证了根节点 \texttt{A[1]} 是整个堆中的最大值。

\subsection{C 代码 — 堆性质维护 Maxheapify}

排序的核心是一个叫 \texttt{Maxheapify} 的函数：
当下述三个节点中有节点违反最大堆性质时，
把最大的元素\underline{换到}父节点位置，然后递归下去。

\lstinputlisting[style=cstyle, firstline=18, lastline=38,
    caption=algo/src/sort/heap\_sort.c — Maxheapify]
{../algo/src/sort/heap_sort.c}

它的工作方式可以概括为：

\begin{enumerate}
    \item 看当前节点 \texttt{i} 和它两个孩子里谁最大
    \item 如果当前节点已经最大，什么也不做
    \item 否则把当前节点和最大的那个孩子交换
    \item 到被交换的那一侧递归调用 \texttt{Maxheapify}
\end{enumerate}

因为每次最多下\emph{一层}，所以一棵有 \(n\) 个节点的树，
深度最多 \(O(\log n)\)，每次调用 \texttt{Maxheapify} 的时间为 \(O(\log n)\)。

\subsection{C 代码 — 建堆 BuildmaxHeap}

建堆从\textbf{最后一个非叶节点}开始（也就是 \texttt{Parent(n)}），
从后往前对每个节点调用 \texttt{Maxheapify}。
这样保证在处理第 \texttt{i} 层时，\texttt{i+1} 层已经是堆。

\lstinputlisting[style=cstyle, firstline=40, lastline=47,
    caption=algo/src/sort/heap\_sort.c — BuildmaxHeap]
{../algo/src/sort/heap_sort.c}

为什么从 \texttt{Parent(n)} 开始而不是 \texttt{n}？
因为超过 \texttt{Parent(n)} 的节点全都是\underline{叶子}，
一个元素的堆本来就满足堆性质，不需要处理。

\subsection{C 代码 — Heapsort 主函数}

排序只需要：
1. 把原数组建成最大堆
2. 反复把根节点（最大值）换到数组末尾，
   堆大小减 1，再从根节点维护堆性质

\lstinputlisting[style=cstyle, firstline=48, lastline=60,
    caption=algo/src/sort/heap\_sort.c — Heapsort]
{../algo/src/sort/heap_sort.c}

第 54--55 行的 \texttt{temp = A[1]; A[1] = A[i]; A[i] = temp;}
就是\textbf{把根换到末尾}的操作，相当于
把当前的最大值放到了它最终的有序位置。

\subsection{C 代码 — 辅助函数与测试}

三个索引计算函数很简单，但它们让排序逻辑语义清晰：

\lstinputlisting[style=cstyle, firstline=1, lastline=17,
    caption=algo/src/sort/heap\_sort.c — Parent/Leftchild/Rightchild]
{../algo/src/sort/heap_sort.c}

\subsection{C 代码 — main 测试程序}

\lstinputlisting[style=cstyle, firstline=62, lastline=73,
    caption=algo/src/sort/heap\_sort.c — main 测试]
{../algo/src/sort/heap_sort.c}

测试数组 \texttt{\{0,5,2,4,6,1,3,7,8\}} 中
\texttt{A[0]=0} 被忽略（占位），对 \texttt{A[1..8]} 排序。
注意 \texttt{sizeof(A)/sizeof(A[0]) - 1} 正好跳过了占位元素，
得到有效长度 \texttt{n=8}。

\subsection{时间复杂度}

| 操作 | 一次调用 | 总次数 |
|------|---------:|-------:|
| \texttt{Maxheapify} | \(O(\log n)\) | \(O(n)\) 次在建堆阶段 + \(O(n)\) 次排序阶段 |
| \texttt{BuildmaxHeap} | — | \(O(n)\)（紧界，不是 \(O(n \log n)\)！） |
| \texttt{Heapsort} | — | \(O(n \log n)\) |

为什么 \texttt{BuildmaxHeap} 是 \(O(n)\) 而不是 \(O(n \log n)\)？
因为大部分节点在树的底部（叶子附近），\texttt{Maxheapify}
在那些节点几乎立刻返回。精确求和：

\[
\sum_{h=0}^{\lfloor \log n \rfloor} \left\lceil \frac{n}{2^{h+1}} \right\rceil O(h)
= O(n)
\]

而排序阶段每次从根维护堆，深度最大，因此是 \(O(n \log n)\)。

\subsection{与其他排序的对比}

\begin{tabular}{lcc}
\hline
算法 & 最坏时间 & 原地？ \\
\hline
插入排序 & \(O(n^2)\) & 是 \\
归并排序 & \(O(n \log n)\) & 否（需要辅助数组） \\
堆排序 & \(O(n \log n)\) & 是 \\
\hline
\end{tabular}

堆排序在\textbf{最坏情况仍保证} \(O(n \log n)\)，
同时只需要 \(O(1)\) 额外空间，这在嵌入式/资源受限场景很有用。
